\documentclass[12pt]{article}
% єдина преамбула -- nmtbase.sty у корені проєкту (знаходиться і з кореня, і з теки теми, і з підтеки)
\IfFileExists{nmtbase.sty}{\usepackage{nmtbase}}{\IfFileExists{../nmtbase.sty}{\usepackage{../nmtbase}}{\usepackage{../../nmtbase}}}
\fancyhead[C]{\small\color{gray!80} База завдань НМТ 2022--2026 \textendash{} Розділ 3}
\fancyhead[R]{\small\color{mainGreen}\textbf{\href{https://t.me/pvtr2525}{@pvtr2525}}}
\fancyfoot[L]{\small\color{gray!80}\href{https://t.me/pvtr2525}{ТГ: @pvtr2525}}
\fancyfoot[C]{\small\color{gray!80}\thepage}
\fancyfoot[R]{\small\color{gray!80}НМТ 2022--2026}
\usepackage{docmute}
% у зведених документах: розділ = section, тема = subsection, рік = subsubsection
\setcounter{tocdepth}{3}
\renewcommand{\chapterTitle}[1]{\clearpage\phantomsection\addcontentsline{toc}{subsection}{#1}\sectionTitle{#1}}
\renewcommand{\typeTitle}[1]{\needspace{6\baselineskip}\phantomsection\addcontentsline{toc}{subsubsection}{\quad #1}%
\par\vspace{0.3cm}\noindent{\large\bfseries\color{yearOrange}#1}\par\vspace{0.08cm}%
\noindent{\color{yearOrange}\rule{\linewidth}{0.4pt}}\par\nopagebreak\vspace{0.2cm}}
\newcommand{\partTitle}[1]{\clearpage\phantomsection\addcontentsline{toc}{section}{#1}%
\vspace*{0.5cm}\begin{tcolorbox}[colback=titleBg, colframe=mainGreen!60, boxrule=0.8pt, arc=8pt, halign=center, fontupper=\LARGE\bfseries\color{mainGreen}]
#1
\end{tcolorbox}\par\vspace{0.4cm}}
\begin{document}
\thispagestyle{empty}
\vspace*{2cm}
\begin{center}
{\Huge\bfseries\color{mainGreen} РОЗДІЛ 3}\\[0.4cm]
{\LARGE\bfseries\color{mainGreen} Функції}\\[0.6cm]
{\Large База завдань НМТ 2022--2026}\\[0.8cm]
{\large Теми: 20, 21, 24, 25, 28b, 30b, 33, 34}\\[1.2cm]
{\large\color{mainGreen}\textbf{\href{https://t.me/pvtr2525}{Телеграм: @pvtr2525}}}
\end{center}
\clearpage
{\hypersetup{linkcolor=black}\tableofcontents}
\partTitle{РОЗДІЛ 3. ФУНКЦІЇ}
\input{"20.\space Функція\space та\space її\space властивості/завдання.tex"}
\input{"21.\space \space Графіки\space елементарних\space функцій,\space їх\space побудова\space графіків\space функцій\space та\space їх\space перетворення/завдання.tex"}
\input{"24.\space Арифметична\space прогресія/завдання.tex"}
\input{"25.\space Геометрична\space прогресія/завдання.tex"}
\input{"28.\space Показникова\space функція.\space Показникові\space рівняння/Показникові\space функція\space і\space вирази/завдання.tex"}
\input{"30.\space Логарифм.\space Логарифмічна\space функція/завдання\space функція.tex"}
\input{"33.\space Похідна/завдання.tex"}
\input{"34.\space Первісна/завдання.tex"}
\end{document}
